\section{引言}
\label{sec:introduction}

驾驶仿真与场景编辑需要能够支持相机移动、车辆运动和人体姿态变化的三维场景~\cite{yang2023unisim,zhou2024hugsim}。然而，驾驶日志中的图像只提供采集视角和当时参与者状态下的可见表面~\cite{sun2020waymo,caesar2020nuscenes}。移动相机、转动车辆或改变行人姿态，都会暴露此前未观测的区域（图~\ref{fig:teaser}）。因此，场景重建需要补全这些区域的几何与外观，并使补全后的车辆和人体表面能够随所属对象运动。

近期前馈方法通过融合多帧图像，直接预测驾驶场景的三维几何与高斯表示~\cite{yang2025storm,chen2025dggt,yu2026streetforward,chi2026pointforward}。运动与可见性估计对齐动态对象，聚合输入序列中可见的表面。新视角或对象运动暴露出整个序列都未曾拍摄的表面，重建需要推断其几何与外观。

生成式重建通过资产先验、高斯补全或生成视图推断未观测内容~\cite{yang2025genassets,liao2025complete,streetcrafter,zhu2025worldsplat}。与此同时，可组合重建将背景与运动对象分开表示，支持车辆移动和人体姿态控制~\cite{yan2024street,chen2025omnire,zhou2024hugsim}。将生成补全用于这类场景，需要使补出的表面随所属对象运动。对象仅被部分观测，背景与其他对象的信息可为推断其缺失内容提供补充上下文。

为此，我们提出 DriveCompose，从驾驶图像序列中构建可独立控制、并经过生成补全的三维场景。我们在世界坐标系中表示背景，在刚体资产自身的坐标系中表示其表面，并在参考姿态下表示人体、保留其骨架绑定~\cite{yan2024street,chen2025omnire}。补全直接在这些表示中进行，使生成的刚体和人体表面能够随所属对象运动或改变姿态（第~\ref{sec:representation} 节）。在此基础上，模型通过共享场景上下文，联合补全背景与各对象的缺失内容。

图~\ref{fig:overview} 概括构建与复用流程。视觉编码恢复几何与运动，规范聚合对齐观测，类型专属编码器生成组件潜码。局部注意力建模组件内部结构，跨组件记忆交换场景上下文，通过流匹配补全缺失内容~\cite{lipman2022flowmatching}。高斯解码器保留刚体与骨架绑定~\cite{kerbl2023gaussians}，使状态变化复用同一份完成内容（第~\ref{sec:composition} 节）。

在 Waymo 的统一相机与参与者状态条件下，DriveCompose 改善了重建与轨迹外渲染。配对合成测试中，新增可见表面优于独立补全；真实编辑在不同参与者状态下复用同一份完成资产（第~\ref{sec:experiments} 节）。

本文的贡献包括：
\begin{itemize}
\item 提出 DriveCompose，从图像序列联合补全背景、车辆与人体，构建在相机、刚体和骨架状态变化时可复用的可组合高斯资产。
\item 设计基于条件流匹配的跨组件潜空间补全模型。通过局部注意力建模各组件内部的空间关系，通过跨组件记忆汇集背景与参与者的特征，为各组件的缺失内容生成提供场景上下文。
\item 通过条件重建、轨迹外渲染、资产编辑与组件消融评估框架及其关键机制。按可见率分组的诊断将联合补全的收益定位到背景与参与者中严重缺失的区域。
\end{itemize}
